\clearpage
\thispagestyle{empty}

\begin{center}
	\vspace*{4cm}
	{\Huge \textbf{KẾT THÚC PHẦN 1}} \\[1.5cm]
	
	\rule{0.8\textwidth}{0.4pt} \\[0.8cm]
	
	\parbox{0.9\textwidth}{
		\centering
		\textit{
			Phần 1 của giáo trình đã cung cấp cho bạn một nền tảng lý thuyết toàn diện, 
			từ những nguyên lý ngôn ngữ học cơ bản, các mô hình thống kê, 
			các kiến trúc nơ-ron kinh điển, cho đến sự thống trị của Transformer 
			và các kỹ thuật tiên tiến để tinh chỉnh, căn chỉnh, và đánh giá các 
			Mô hình Ngôn ngữ Lớn. Với nền tảng này, bạn đã sẵn sàng để bước sang 
			Phần 2, nơi chúng ta sẽ biến những lý thuyết này thành các kỹ năng và 
			quy trình thực chiến để xây dựng các ứng dụng NLP thực tế.
		}
	} \\[1cm]
	
	\rule{0.8\textwidth}{0.4pt} \\[2cm]
\end{center}

\noindent
\textbf{Tóm lược nội dung Phần 1:}
\begin{itemize}
	\item \textbf{Chương 1 – Nhập môn và Nguyên lý nền tảng:} 
	Định nghĩa NLP, mối quan hệ với AI, các kỷ nguyên phát triển, 
	kiến thức ngôn ngữ học, toán học, và các vấn đề đạo đức trong NLP.
	\item \textbf{Chương 2 – Biểu diễn văn bản và Mô hình thống kê:} 
	Các kỹ thuật BoW, TF-IDF, N-gram, smoothing, 
	mô hình phân loại, học chuỗi, và LDA.
	\item \textbf{Chương 3 – Các kiến trúc mạng nơ-ron kinh điển:} 
	Word embeddings (Word2Vec, GloVe, FastText), 
	autoencoder, RNN/LSTM/GRU, seq2seq với attention, 
	CNN cho NLP, contextualized embeddings, và các mô hình sinh pre-Transformer.
	\item \textbf{Chương 4 – Kỷ nguyên Transformer và LLMs:} 
	Kiến trúc Transformer gốc, phân loại LLMs, tokenization hiện đại, 
	biến thể cho ngữ cảnh dài, MoE.
	\item \textbf{Chương 5 – Tinh chỉnh và Căn chỉnh mô hình:} 
	Fine-tuning, PEFT (LoRA, QLoRA, Adapter), prompt engineering, 
	instruction tuning, RLHF và các kỹ thuật alignment khác.
	\item \textbf{Chương 6 – Hệ thống và Kiến trúc nâng cao:} 
	RAG, tối ưu hóa mô hình (distillation, quantization, pruning), 
	mô hình đa phương thức, hệ thống tác tử, world models, 
	và biểu diễn nâng cao cho truy xuất.
	\item \textbf{Chương 7 – Đánh giá và Benchmark:} 
	Các metric kinh điển (perplexity, F1, BLEU, ROUGE), 
	benchmark (GLUE, SuperGLUE, đánh giá emergent abilities), 
	thách thức trong đánh giá, và xu hướng “LLM-as-a-Judge”.
\end{itemize}